\scanpage{14}

% Clipped upper-left edge ink, source x=23--50, y=117--174.
% Its textual status and orientation are uncertain; no lost letters are supplied.
\margin{\figure{\unsure{\begin{tikzpicture}[x=.22mm,y=-.22mm,line width=.45pt,line cap=round,line join=round]
\draw (0,0)--(5,4) .. controls (15,2) and (22,8) .. (17,13) .. controls (12,18) and (6,8) .. (0,3);
\draw (0,24) .. controls (11,31) and (25,37) .. (25,29) .. controls (25,25) and (20,25) .. (18,25);
\draw (1,49) .. controls (4,51) and (3,54) .. (1,56);
\end{tikzpicture}}}}

b) $\forall\,F$ injectif sur $\unsure{Y}$, $\unsure{q|g}^{*}(F)$ est injectif et $F\unsure{\mathord{-}|\to}q_{*}q^{*}F$ est un épimorphisme (ex: $q$ inclusion d'un ouvert $U\hookrightarrow e_{Y}$).

On prend
\[
\mathbb R\Gamma_{Y}\bigl(\operatorname{Ker}\bigl(\underbrace{C(F}_{\text{rés. inj.}})\to q_{*}q^{*}\unsure{(|C}C(F)\bigr)\bigr)\bigr)
\]
% Right margin beside the first display: three faint, detached traces.
\margin{\figure{\unsure{\begin{tikzpicture}[x=.5mm,y=-.5mm,line width=.3pt,line cap=round]
\draw[gray] (0,0)--(1.2,0);
\fill (0.7,1.7) ellipse (.35 and .45);
\draw[gray] (.7,2)--(.2,3.7);
\fill[gray] (0,5) ellipse (.3 and .65);
\end{tikzpicture}}}}

Dans le cas général, la difficulté provient du fait que le cône d'un morphisme de complexes (tel que
% Clipped left edge, level with the second display.
\margin{\figure{\unsure{\begin{tikzpicture}[x=.24mm,y=-.24mm,line width=.45pt,line cap=round]
\draw (0,0)--(.5,16);
\draw (0,20)--(12,18.5);
\fill (0,32) ellipse (.8 and 1.3);
\end{tikzpicture}}}}
\[
F\to q_{*}\deleted{\unsure{C|(}}\bigl(q^{*}(F)\deleted{\unsure{)|C}}\bigr)\deleted{\unsure{*|q_{*}}}\qquad\text{)}
\]
n'est pas fonctoriel (dans la catégorie dérivée) p. r. à la flèche dont on veut prendre le cône.
% Detached left-margin points at source (23,813) and (22,914).
\margin{\unsure{.|\phantom{.}}}
\margin{\unsure{.|\phantom{.}}}

% Clipped left-edge curve beside the Et pourtant paragraph, source y=1040--1059.
\margin{\figure{\unsure{\begin{tikzpicture}[x=.35mm,y=-.35mm,line width=.45pt,line cap=round]
\draw (0,0) .. controls (8,5) and (8,11) .. (0,13);
\end{tikzpicture}}}}
Et pourtant, dans le cas particulier \unsure{actuel,|actuel|actuels,} il devrait y avoir un choix fonctoriel. ESt$\underset{\unsure{\cdot|\phantom{\cdot}}\mkern12mu}{\text{-}}$ce \unsure{évident|evident} ?
% The point under the hyphen is spatial only, not a subscript or an asserted punctuation role.
% Separate short high stroke in the right margin, source (1069,1154).
\margin{\unsure{'|\phantom{'}}}
% Clipped left-edge curve and low point, source y=1200--1235.
\margin{\figure{\unsure{\begin{tikzpicture}[x=.35mm,y=-.35mm,line width=.45pt,line cap=round]
\draw (0,0) .. controls (8,6) and (8,13) .. (0,16);
\fill (0,33) ellipse (.55 and .8);
\end{tikzpicture}}}}
% Isolated point in the gap above the rubric, source (510,1240).
\margin{\unsure{.|\phantom{.}}}

% Two tiny clipped boundary marks at the left edge, level with the rubric's underline.
\margin{\figure{\unsure{\begin{tikzpicture}[x=.4mm,y=-.4mm,line width=.5pt,line cap=round]
\draw (0,0)--(2,1.2);
\fill (1,12) circle (.8);
\end{tikzpicture}}}}
\uline{Question pour Illusie} : Dans sa \unsure{théorie|theorie} des déform.tions de schémas en \unsure{groupes|groupe} plats, il tombe sur des $H^{3}(B_{G}/X,\quad)$ resp. des $\operatorname{Ext}^{2}(\unsure{K;,|\underline{K}.,|X;,|\underline{X}.,}\quad)$. Peut-on courtcircuiter sa théorie via la théorie (supposée écrite) des Gr-champs - resp. via la théorie des champs de Picard ? J'ai

% Clipped left-edge blob and low point beside the final paragraph, source y=1410--1434.
\margin{\figure{\unsure{\begin{tikzpicture}[x=.4mm,y=-.4mm,line width=.5pt,line cap=round]
\draw (0,0) .. controls (4,4) and (4,8) .. (0,10);
\fill (0,9) ellipse (1.3 and 1.5);
\fill (0,21) ellipse (.5 and .75);
\end{tikzpicture}}}}
% Bottom margin near the page centre: one larger dark point and a smaller point.
\margin{$\unsure{\bullet\,\cdot|\phantom{\bullet\,\cdot}}$}
% archive sheet number: 13
